\documentclass[11pt]{article}
\usepackage[margin=1in]{geometry}
\usepackage{amsmath,amssymb,amsthm,mathtools}
\IfFileExists{lmodern.sty}{\usepackage{lmodern}}{}
\usepackage[expansion=false]{microtype}
\usepackage{booktabs,array}
\usepackage{enumitem}
\usepackage{xcolor}
\usepackage[colorlinks=true,linkcolor=blue!50!black,citecolor=green!35!black,urlcolor=blue!50!black]{hyperref}
\usepackage[nameinlink]{cleveref}

\newtheorem{theorem}{Theorem}[section]
\newtheorem{proposition}[theorem]{Proposition}
\newtheorem{lemma}[theorem]{Lemma}
\newtheorem{corollary}[theorem]{Corollary}
\theoremstyle{definition}
\newtheorem{definition}[theorem]{Definition}
\newtheorem{example}[theorem]{Example}
\newtheorem{question}[theorem]{Question}
\theoremstyle{remark}
\newtheorem{remark}[theorem]{Remark}
\newtheorem*{mainthm}{Main results}

\newcommand{\D}{\mathsf{D}}
\newcommand{\R}{\mathbb{R}}
\newcommand{\Cpl}{\operatorname{Cpl}}
\newcommand{\diag}{\operatorname{diag}}
\newcommand{\TV}{\mathrm{TV}}
\newcommand{\T}{\mathsf T}
\newcommand{\Hup}{\overline{H}}
\newcommand{\Hlo}{\underline{H}}

\title{Bands:\\
Perturbation, Tearing and Breaking of Coupling Structures,\\ and the Entropy of Changing Probabilities}
\author{Vinay K. Awasthi\\
Department of Applied Mathematics and Statistics\\
Johns Hopkins University\thanks{\texttt{check\_bands.py} (17 checks) and \texttt{check\_dynamics.py} (9 checks) verify numerically the statements proved here, the synthetic
validation and every quoted number; results cited from the other papers are checked there. The pilots are reproduced by the
scripts in the folder \texttt{pilot}.}}
\date{October 8, 2026, draft 2}

\begin{document}
\maketitle

\begin{abstract}
Probabilities in practice are not fixed. They are revised as information arrives, and revised again when the information
turns out to be unreliable. We model the state of knowledge as a \emph{band}, a closed convex set of distributions, typically
given by elementwise limits $(\mathrm{lo},\mathrm{hi})$, and ask how far the algebraic and geometric structures of couplings
tolerate perturbations of a band before they change. There are three regimes. Inside a chamber of margins the structures
\emph{bend}: the gluing algebra keeps its type, the coupling polytope keeps its combinatorial type, and table counts stay
polynomial. On a wall they \emph{tear}: the polytope degenerates, its vertex supports change, and for two-row tables the
counts jump by explicit terms. At the boundary of the
simplex they \emph{break}: an outcome loses its support and the algebra drops dimension. The total-variation distance to
tearing is $m_T=\min\lvert f(I)-g(J)\rvert$ and to breaking is $m_B=\min_y f_y$, both exact. Entropy is continuous in all three
regimes, and the regimes are told apart by its derivatives: the product coupling's entropy is smooth everywhere inside
the simplex, a vertex coupling whose cell vanishes on a wall has a $d\log(1/d)$ cusp there, and at the boundary the gradient
diverges. On bands, entropy becomes an interval whose upper end is a water-filling and whose lower end sits at a vertex.
Information acts by a monoid of updates on bands, not a group: conditioning commutes and is idempotent, Jeffrey updates do not
commute, discounting widens bands and noisy kernels contract them. The ratio of the margin to the speed of an update is a
tear time, which lower-bounds the number of updates before a structure tears. Assembling the fibres of compatible laws into families
over a band gives one polytope with a computable entropy interval; limits become compatibility (cones are compatible laws) and colimits
assembly of fibres; and the
parameter space splits into structural shards. Along explicit paths each structure is seen to bend, tear or break. A break has
second-order effects: global compatibility can be lost only through a break of the maximum-entropy filler, with an early warning
that is sharp in an example; a zero at level one of the tower removes a model for ever, while reliability at level three bounds how
fast the calibrated view can change and prevents the break. Two pilots on public data, NHANES subgroups and
COVID-19 ensemble forecasts with their revision history, show much of the predicted behaviour: bands shrink up an
information ladder, narrow structural bands must be combined with sampling bands to predict new data, decisive bands cross
walls by straddling them, small margins predict tearing, and forecasts built on data that was later heavily revised lose
almost half their coverage. The speed term of the tear time does not yet improve on the margin. A protocol for clinical data separates
the statements that are theorems from those that can fail, and it is validated on synthetic data with planted effects.
\end{abstract}

\tableofcontents

%=====================================================================
\section{Introduction}\label{sec:intro}
%=====================================================================

The papers of this programme treat a distribution as given: a margin $f$, a coupling of $f$ with $g$, a posterior at some
level of a tower. In practice the distribution is an estimate. It moves when new data arrive, and it moves again when old data
are corrected. A forecast of a week's deaths is made four weeks ahead and revised three times before the week ends, and the counts it is scored
against are themselves revised for weeks afterwards. A subgroup risk is known only within limits, and those limits narrow
as more cross-tabulations are published.

This paper asks three questions. How large a change of the distributions can each structure absorb before it changes?
What does entropy, as a function on sets of distributions, do as they move? And what is the right notion of the rate at
which information changes them? The setting is a \emph{band}: a closed convex set of distributions, in the simplest case all
$p$ with $\mathrm{lo}\le p\le\mathrm{hi}$ (\cref{sec:bands}). A band is what the second level of the posterior tower
\cite{Tower} looks like from the first, and reliability is a third level: discounting a law with reliability $\alpha$ gives
exactly the band $[\alpha p,\alpha p+1-\alpha]$ (\cref{prop:discount}).

\begin{mainthm}
\leavevmode
\begin{enumerate}[label=(\Alph*),leftmargin=*]
\item \emph{Three regimes, with exact distances} (\cref{thm:distances,thm:regimes}). Structures bend inside a chamber,
  tear on a wall and break at the boundary of the simplex; chambers are taken inside the open simplices. The total-variation distance from a pair of margins to the
  nearest wall is $m_T=\min_{I,J}\lvert f(I)-g(J)\rvert$, and to the boundary is $m_B=\min_yf_y$.
\item \emph{Tolerance of a band} (\cref{prop:bandsafe}). For elementwise bands the extremes of every event probability
  have closed forms, so whether a band can tear is decided by finitely many sign checks.
\item \emph{Entropy across the regimes} (\cref{prop:entropyregimes}). Entropy is continuous throughout. The product
  coupling's entropy is smooth everywhere inside the simplex; the entropy of a vertex coupling has a $d\log(1/d)$ cusp
  where one of its cells vanishes on a wall; at the boundary the gradient diverges. So tearing is visible to extremal
  structures but not to the product coupling, and breaking to both.
\item \emph{Entropy of a band} (\cref{prop:upper,prop:lower,prop:entropyspeed,prop:level2}). The upper entropy is a
  water-filling, the lower entropy sits at a vertex, both move by at most $t\log(n-1)+h(t)$ when the band moves by
  $t\le1-1/n$ in total variation, and the entropy of a mixture splits into expected entropy plus the information about which model holds.
  The entropy gap between the conditionally independent gluing and a joint with the same margins is the conditional
  mutual information.
\item \emph{The update monoid} (\cref{prop:monoid,prop:dilation,prop:loop}). Updates of bands form a monoid. Conditioning on events
  commutes and is idempotent; Jeffrey updates do not commute, but they preserve odds ratios, so a loop of them leaves a residue
  that iterated fitting removes; discounting widens a band; a kernel contracts it by its Dobrushin coefficient; and
  conditioning can dilate a band from a point to $[0,1]$.
\item \emph{Tear time} (\cref{prop:teartime}). If each update moves the margins by at most $v$ in total variation, no
  structure tears before $\lceil m_T/v\rceil$ updates.
\item \emph{Families, limits, shards} (\cref{prop:family,def:shards,prop:updateregimes}). The family of fibres over a polyhedral band is
  a polytope whose upper entropy is the sum of the margins' upper entropies; cones are compatible laws; breaking is created only by
  conditioning, Jeffrey updates with a zero target and kernels with a zero column, and removed by kernels with positive entries.
\item \emph{Stability along paths} (\cref{ex:path22,ex:path3}). Structure by structure, what bends, tears and breaks along a $2\times2$ path
  through two walls and along the three-bit path.
\item \emph{Second-order breaking} (\cref{thm:lossbybreak,prop:tolerance}). Compatibility is lost only through a break of the filler,
  and the filler's break margin divided by $\lVert M^+\rVert$ is a tolerance, attained on part of the three-bit path.
\item \emph{The tower under evidence} (\cref{ex:cromwell,prop:morph}). A level-one zero removes a model for ever; linear discounting with reliability $\alpha$
  bounds the odds change per observation by $n/(1-\alpha)$ and prevents the break.
\item \emph{Evidence} (\cref{sec:pilots,sec:protocol}). Two pilots on public data, and a test protocol for clinical
  data in which theorems and falsifiable claims are kept apart, validated on synthetic data.
\end{enumerate}
\end{mainthm}

\paragraph{What is not claimed.}
\begin{itemize}[leftmargin=*,nosep]
\item The dynamics of information is not a Galois action. Updates do not invert, so they form a monoid; its invertible
  elements are relabellings, which supply symmetries and invariants but not the motion (\cref{rem:galois}).
\item A loop of updates on margins does not exhibit a topological holonomy. It leaves a residue, but the updates
  preserve the interaction parameters, and further fitting to the original margins removes the residue (\cref{prop:loop});
  genuine holonomy needs updates that change interactions, and is left open.
\item In the forecast pilot an analogue of the tear time predicted tearing, but no better than the margin alone
  (\cref{sec:fhub}). Its speed term does not yet earn its place.
\item Limits are not claimed to exist as universal cones: the category of finite probability spaces has no products, and we use
  cones, compatible laws, with entropy choosing a canonical one. Our shards are unrelated to the shards of hyperplane arrangements.
\item In-sample containment of the true value in the NHANES bands is a sanity check, not evidence; the out-of-sample
  version is in \cref{sec:nhanes}.
\item In data drawn from one classical population, a common joint distribution always exists. Incompatibility of observed
  margins therefore measures selection between contexts, not contextuality of the kind in \cite{Conditionals}
  (\cref{sec:protocol}).
\end{itemize}

\paragraph{Relation to the other papers.} The walls and chambers are those of \cite[Def.~6.3, Prop.~6.6]{Conditionals}, and the
chamber-identification bound is \cite[Prop.~12.2]{Conditionals}. The gluing algebra and the table counts across walls are
\cite[Thm.~2.2, Prop.~7.3]{Algebras}. The band as a level-two shadow and the Dobrushin contraction come from
\cite[Prop.~9.14, Prop.~11.2]{Tower}. The precision ladder of \cite[Thm.~3.5]{Attribution} reappears as the information
ladder of the NHANES pilot. The fixed-fibre geometry, constraint branes and torus branes, is in the companion note
\cite{Bridges}; this paper supplies its dynamics. The parity horns and the signed-measure gluing are from \cite[Thm.~4.3]{Conditionals} and
\cite[Prop.~35.1]{ThreeLayers}. Each of these papers can later gain a short subsection citing the part of this one it
needs.

%=====================================================================
\section{Bands}\label{sec:bands}
%=====================================================================

Let $Y$ be a finite set with $n$ elements and $\D(Y)$ the simplex of distributions on it.

\begin{definition}\label{def:band}
A \emph{band} on $Y$ is a nonempty closed convex subset $B\subseteq\D(Y)$. An \emph{elementwise band} is
$B(\mathrm{lo},\mathrm{hi})=\{p\in\D(Y):\mathrm{lo}\le p\le\mathrm{hi}\}$ for vectors $0\le\mathrm{lo}\le\mathrm{hi}\le1$; it is
nonempty if and only if $\sum\mathrm{lo}\le1\le\sum\mathrm{hi}$. Its \emph{tightening} is
\[
\mathrm{lo}'_y=\max\Bigl(\mathrm{lo}_y,\ 1-\sum_{z\ne y}\mathrm{hi}_z\Bigr),\qquad
\mathrm{hi}'_y=\min\Bigl(\mathrm{hi}_y,\ 1-\sum_{z\ne y}\mathrm{lo}_z\Bigr).
\]
\end{definition}

Elementwise bands are the probability intervals of de Campos, Huete and Moral \cite{deCampos1994}, and general bands are
credal sets in the sense of Walley \cite{Walley1991}.

\begin{proposition}[Tightening and event probabilities]\label{prop:tighten}
Let $B=B(\mathrm{lo},\mathrm{hi})$ be nonempty. Then $B(\mathrm{lo}',\mathrm{hi}')=B$, every tightened bound is attained by
some $p\in B$, and for every event $I\subseteq Y$,
\[
\min_{p\in B}p(I)=\max\Bigl(\sum_{y\in I}\mathrm{lo}_y,\ 1-\sum_{y\notin I}\mathrm{hi}_y\Bigr),\qquad
\max_{p\in B}p(I)=\min\Bigl(\sum_{y\in I}\mathrm{hi}_y,\ 1-\sum_{y\notin I}\mathrm{lo}_y\Bigr).
\]
\end{proposition}
\begin{proof}
For $p\in B$, $p(I)=1-p(I^c)\ge1-\sum_{I^c}\mathrm{hi}$ and $p(I)\ge\sum_I\mathrm{lo}$, so the stated minimum is a lower bound. It
is attained by setting $p=\mathrm{hi}$ on $I^c$ and $p=\mathrm{lo}$ on $I$ and then distributing the remaining mass
$1-\sum_I\mathrm{lo}-\sum_{I^c}\mathrm{hi}$, which has the sign that makes this possible, on $I$ if it is positive and on
$I^c$ if it is negative, within the bounds; this is a greedy fill, possible because $\sum\mathrm{lo}\le1\le\sum\mathrm{hi}$. The
maximum is symmetric. Singletons $I=\{y\}$ give the tightened bounds, which are therefore attained, and they cut out the same
set.
\end{proof}

\begin{proposition}[Reliability gives a band]\label{prop:discount}
Let $p\in\D(Y)$ and $\alpha\in[0,1]$. The set $\{\alpha p+(1-\alpha)q:q\in\D(Y)\}$ of laws obtained by trusting $p$ with
reliability $\alpha$ and allowing anything otherwise, the credal set of Shafer's discounting \cite{Shafer1976}, is exactly
the elementwise band $B(\alpha p,\ \alpha p+(1-\alpha)\mathbf 1)$.
\end{proposition}
\begin{proof}
Each $\alpha p+(1-\alpha)q$ lies between $\alpha p$ and $\alpha p+(1-\alpha)$ coordinatewise. Conversely, if $r$ lies in
the band and sums to $1$, then $q=(r-\alpha p)/(1-\alpha)$ is nonnegative and sums to $1$ (for $\alpha<1$; for $\alpha=1$ both
sets are $\{p\}$).
\end{proof}

\paragraph{Bands from the tower.} A finitely supported law $\Pi=\sum_kw_k\delta_{\rho_k}$ on $\D(Y)$, a level-two element of the
tower \cite{Tower}, generates the band $\operatorname{conv}\{\rho_k\}$, which contains the barycenter $\sum_kw_k\rho_k$. The
band forgets the weights $w_k$, and that is the precise sense in which it is the level-one shadow of level two. A level-three
law, a law on such $\Pi$, models doubt about the reliability of the level-two assessment; \cref{prop:discount} is the
simplest instance, with $\Pi$ supported on $\{p\}$ and the discounting weight standing for the doubt.

\paragraph{Banded margins.} When both margins are only known to lie in bands, the cells of their couplings are bounded by the
Fréchet bounds taken at the extremes of the bands.

\begin{proposition}[Fréchet bounds over bands]\label{prop:frechet}
Let $F=B(\mathrm{lo}^f,\mathrm{hi}^f)$ and $G=B(\mathrm{lo}^g,\mathrm{hi}^g)$ be tightened. Over all $f\in F$, $g\in G$ and
$\alpha\in\Cpl(f,g)$,
\[
\max\bigl(0,\ \mathrm{lo}^f_i+\mathrm{lo}^g_j-1\bigr)\ \le\ \alpha_{ij}\ \le\ \min\bigl(\mathrm{hi}^f_i,\ \mathrm{hi}^g_j\bigr),
\]
and both bounds are attained.
\end{proposition}
\begin{proof}
For fixed margins the Fréchet bounds $\max(0,f_i+g_j-1)\le\alpha_{ij}\le\min(f_i,g_j)$ are sharp. They are monotone in
$f_i$ and $g_j$, and the two bands are independent, so the extremes are attained at $f_i,g_j$ both minimal or both maximal,
which is possible by \cref{prop:tighten}.
\end{proof}

%=====================================================================
\section{Three regimes}\label{sec:regimes}
%=====================================================================

Let $f$ and $g$ be distributions on sets with $n$ and $m$ elements. A \emph{wall} is a hyperplane $f(I)=g(J)$ with $I,J$
nonempty and proper, and the \emph{chambers} are the components of the complement of the walls in the product of the open
simplices \cite[Def.~6.3]{Conditionals}. Off the open simplices the polytope can degenerate without any wall, which is the
breaking regime.

\begin{definition}\label{def:margins}
The \emph{tear margin} of $(f,g)$ is $m_T(f,g)=\min_{I,J}\lvert f(I)-g(J)\rvert$, over nonempty proper $I,J$. The \emph{break
margin} of $f$ is $m_B(f)=\min_yf_y$, and of the pair it is $\min(m_B(f),m_B(g))$. Distances between pairs are measured by
$\TV(f,f')+\TV(g,g')$, with $\TV(p,q)=\frac12\sum\lvert p-q\rvert$.
\end{definition}

\begin{theorem}[Exact distances]\label{thm:distances}
\leavevmode
\begin{enumerate}[label=(\roman*)]
\item The distance from $(f,g)$ to the union of the walls is $m_T(f,g)$.
\item The distance from $f$ to the face $\{p_y=0\}$ of the simplex is $f_y$, so the distance to its boundary is
  $m_B(f)$.
\end{enumerate}
\end{theorem}
\begin{proof}
(i) For any event, $\lvert f(I)-f'(I)\rvert\le\TV(f,f')$, so $f(I)-g(J)$ changes by at most $\TV(f,f')+\TV(g,g')$ and no
pair closer than $m_T$ lies on a wall. Conversely, if $f(I)-g(J)=d>0$, moving mass $d$ from $I$ to $I^c$ in $f$ costs exactly
$d$ in total variation and lands on the wall; this is possible because $f(I)\ge d$ and $I^c\ne\emptyset$. (ii) Moving the
mass $f_y$ elsewhere costs $f_y$, and any $p$ with $p_y=0$ has $\TV(f,p)\ge f_y$.
\end{proof}

By \cite[Prop.~12.2]{Conditionals}, $N\ge2\log(2W/\eta)/m_T^2$ samples identify the chamber with probability at least
$1-\eta$, where $W$ is the number of walls. So the tear margin is at once a geometric distance and a sample size.

\begin{theorem}[Bending, tearing, breaking]\label{thm:regimes}
\begin{enumerate}[label=(\roman*)]
\item \emph{Bending.} On the open simplex the gluing algebra $A_f$ is isomorphic to $\R\times M_{n-1}(\R)$ by an isomorphism
  that depends continuously on $f$ \cite[Thm.~2.2]{Algebras}. On a chamber the coupling polytope $\Cpl(f,g)$ is simple and
  of constant combinatorial type \cite[Thm.~6.4, Prop.~6.6]{Conditionals}; its vertices have constant supports, the
  spanning trees of the complete bipartite graph that are feasible, each vertex is an affine function of $(f,g)$ with exactly
  $n+m-1$ positive cells, and for two-row tables the number of integer tables is a single polynomial \cite[Prop.~7.3]{Algebras}.
\item \emph{Tearing.} On a wall the coupling polytope is degenerate: some vertex has fewer than $n+m-1$ positive cells
  \cite{DeLoeraKim2014}. Crossing the wall changes the set of feasible spanning trees, and with it the labelled structure of
  the vertices; the number of vertices may or may not change (for $2\times2$ the polytope is a segment on both sides). For
  $2\times m$ tables the change in the count of integer tables is one binomial term, and a corner cut of lattice depth $d$
  removes $\binom{d+1}2$ tables \cite[Prop.~7.3]{Algebras}; in the $2\times3$ case of \cite[Prop.~10.6]{Conditionals} the area
  changes by $\frac12s^2$.
\item \emph{Breaking.} When $f_y\to0$, the gluing product $a\cdot b=a\,\diag(f)^{-1}b$ involves $\diag(f)^{-1}$ of norm $1/m_B$,
  and the Fisher information of the multinomial family, $\diag(f)^{-1}$ on the tangent space, diverges. On the face
  $f_y=0$ the algebra is replaced by $A_{f\vert_S}$ for the support $S$, of dimension $(\lvert S\rvert-1)^2+1$, strictly smaller.
\end{enumerate}
\end{theorem}
\begin{proof}
(i) The isomorphism is the congruence $a\mapsto\diag(f)^{-1/2}a\,\diag(f)^{-1/2}$. For nondegenerate margins the transportation
polytope is simple, and its vertices are the basic feasible solutions, whose supports are spanning trees of the complete bipartite
graph. Removing an edge of a tree splits its nodes into a part containing rows $I$ and columns $J$ and the rest, and the flow on
that edge is $\pm(f(I)-g(J))$, the sign depending on whether the edge leaves the part from a row or from a column. So which trees are
feasible is decided by the signs of the wall functions, which are constant on a chamber, and each basic solution solves a fixed
linear system in $(f,g)$. (ii) The same computation shows that on a wall some edge of some feasible tree carries flow $0$, so the
corresponding vertex is degenerate; that degeneracy occurs exactly on the walls is classical \cite{DeLoeraKim2014}. The jumps are the
cited results. (iii) Direct; the dimension formula is \cite[Thm.~2.2]{Algebras} applied to the support.
\end{proof}

\begin{proposition}[Tolerance of a band]\label{prop:bandsafe}
Let $F$ and $G$ be elementwise bands for $f$ and $g$. The pair of bands meets no wall if and only if, for every wall $(I,J)$,
either $\min_Ff(I)>\max_Gg(J)$ or $\max_Ff(I)<\min_Gg(J)$, with the extremes given by \cref{prop:tighten}. It meets no face of
the simplex if and only if all tightened lower bounds are positive.
\end{proposition}
\begin{proof}
The bands are independent, so the range of $f(I)-g(J)$ over $F\times G$ is the interval between $\min f(I)-\max g(J)$ and
$\max f(I)-\min g(J)$, and it avoids $0$ exactly under the stated condition. The face statement is \cref{prop:tighten} for
singletons.
\end{proof}

So the tolerance of a band is decided by $2(2^{n-1}-1)(2^m-2)$ closed-form numbers, one minimum and one maximum per wall,
and the margin by which it is safe is the smallest gap among them.

%=====================================================================
\section{Entropy on bands and across the regimes}\label{sec:entropy}
%=====================================================================

Write $H(p)=-\sum p_y\log p_y$ and $h(t)=H(t,1-t)$.

\begin{definition}\label{def:bandentropy}
The \emph{upper} and \emph{lower entropy} of a band $B$ are $\Hup(B)=\max_BH$ and $\Hlo(B)=\min_BH$; their difference is the
entropy that is due to imprecision rather than to randomness.
\end{definition}

\begin{proposition}[Upper entropy is a water-filling]\label{prop:upper}
For an elementwise band, $\Hup$ is attained at the unique point $p_y=\operatorname{clip}(c,\mathrm{lo}_y,\mathrm{hi}_y)$, where
$c$ is the level at which these values sum to $1$.
\end{proposition}
\begin{proof}
$H$ is strictly concave, so the maximizer is unique and characterized by the Karush--Kuhn--Tucker conditions
$-\log p_y-1-\lambda+\mu_y-\nu_y=0$ with $\mu_y\ge0$ active only at $\mathrm{lo}_y$ and $\nu_y\ge0$ active only at
$\mathrm{hi}_y$. So $p_y=e^{-1-\lambda}$ wherever no bound is active, $p_y=\mathrm{lo}_y$ only if $e^{-1-\lambda}\le\mathrm{lo}_y$, and
$p_y=\mathrm{hi}_y$ only if $e^{-1-\lambda}\ge\mathrm{hi}_y$; with $c=e^{-1-\lambda}$ this is the clip. Its sum is continuous and
nondecreasing in $c$, so $c$ exists. Algorithms of this kind for credal sets are given by Abellán and Moral \cite{AbellanMoral2003}.
\end{proof}

\begin{proposition}[Lower entropy sits at a vertex]\label{prop:lower}
$\Hlo$ of a polytope band is attained at a vertex. The vertices of an elementwise band are the points with every coordinate but
at most one at a bound, so there are at most $n2^{n-1}$ of them.
\end{proposition}
\begin{proof}
A concave function attains its minimum over a polytope at a vertex. A point of $B(\mathrm{lo},\mathrm{hi})$ with two coordinates
strictly inside their bounds lies on a segment in $B$ obtained by trading mass between them, so it is not a vertex.
\end{proof}

\begin{proposition}[Entropy moves slowly]\label{prop:entropyspeed}
If $\TV(p,q)=t\le1-1/n$, then $\lvert H(p)-H(q)\rvert\le t\log(n-1)+h(t)$, with equality for $p=\delta_y$ and
$q=(1-t,\frac t{n-1},\dots)$ \cite{Zhang2007,Audenaert2007}. Consequently, if two bands are within Hausdorff distance $t$ in total
variation with $t\le1-1/n$, their upper entropies and their lower entropies each differ by at most $t\log(n-1)+h(t)$.
\end{proposition}
\begin{proof}
The first statement is the sharp continuity bound. The bound is increasing in $t$ on $[0,1-1/n]$, since its derivative
$\log(n-1)+\log\frac{1-t}t$ is nonnegative there. For every $p\in B$ there is $p'\in B'$ within $t$, so
$\Hup(B')\ge H(p)-t\log(n-1)-h(t)$; take the maximum over $p$, and argue symmetrically and likewise for minima.
\end{proof}

So an update that moves a band by $t$ in total variation changes its entropy interval by at most $t\log(n-1)+h(t)$: the rate of
change of entropy is controlled by the rate of change of the band.

\begin{proposition}[Entropy across the regimes]\label{prop:entropyregimes}
\begin{enumerate}[label=(\roman*)]
\item The entropy $H(f\otimes g)=H(f)+H(g)$ of the product coupling, the maximum-entropy coupling, is real-analytic on the
  product of the open simplices. It does not see walls.
\item Let $v(f,g)$ be a vertex of $\Cpl(f,g)$ on a chamber, and suppose that, as $(f,g)$ approaches a point of a wall along a
  segment, one of its cells tends to $0$ while the others stay positive. With $t$ the distance along the segment,
  $H(v(t))=H(v(0))+\kappa t\log(1/t)+O(t)$ for some $\kappa>0$, so its derivative in $t$ diverges like $\kappa\log(1/t)$: a cusp.
\item As $f_y\to0$, $\partial H(f)/\partial f_y=-\log f_y-1$ diverges, and so does the corresponding derivative of the
  product-coupling entropy $H(f)+H(g)$.
\end{enumerate}
So entropy is continuous in all three regimes. For a fixed vertex, bending is smooth, tearing is a logarithmic cusp, and
breaking is a divergent gradient; the product coupling sees only the last.
\end{proposition}
\begin{proof}
(i) is immediate. (ii) On the chamber $v$ is affine in $(f,g)$ by \cref{thm:regimes}(i), so the vanishing cell is $\kappa t$ with
$\kappa>0$, and $-\kappa t\log(\kappa t)=-\kappa t\log t+O(t)$; the other cells stay positive, so their terms are smooth. (iii)
Direct.
\end{proof}

\begin{example}
For $2\times2$ couplings with $g=(0.3,0.7)$, the comonotone vertex has $\alpha_{11}=\min(f_1,0.3)$. Its entropy has one-sided slope
$7.55$, $12.16$ and $16.76$ at distances $10^{-3}$, $10^{-5}$ and $10^{-7}$ past the wall $f_1=0.3$, growing by $\log100$ per factor
$100$, while the entropy of the product coupling has equal one-sided slopes there (check B7).
\end{example}

The lower entropy of the coupling polytope, the minimum over its vertices, is a minimum of finitely many analytic functions on
each chamber. It therefore has ordinary kinks inside chambers where the minimizing vertex changes: for $g=(0.3,0.7)$ the
comonotone and countermonotone entropies cross at $f_1=0.5$, inside a chamber. It has a logarithmic cusp at a wall only when a
minimizing vertex degenerates there, and it can be smooth across a wall otherwise. What distinguishes tearing is therefore the
$\log(1/t)$ blow-up of a slope, not non-smoothness as such. The upper end of the entropy interval of the couplings, the product
coupling, is blind to walls altogether.

\begin{proposition}[Entropy of a mixture]\label{prop:level2}
For $\Pi=\sum_kw_k\delta_{\rho_k}$ with barycenter $\bar\rho=\sum_kw_k\rho_k$,
\[
H(\bar\rho)=\sum_kw_kH(\rho_k)+I(K;X),
\]
where $(K,X)$ has law $w_k\rho_k(x)$. The first term is the expected entropy of the outcome under each model, the aleatoric
part, and the second is the information the outcome carries about the model, the epistemic part.
\end{proposition}
\begin{proof}
$H(X)=H(X\mid K)+I(K;X)$.
\end{proof}

\begin{proposition}[The gluing gap]\label{prop:gluinggap}
Let $P$ be a law of $(X,Y,Z)$ and $F=P_{XY}P_{YZ}/P_Y$ the conditionally independent gluing of its margins, the canonical filler
of \cite[Thm.~4.2]{Conditionals}. Then $F$ has the same $(X,Y)$ and $(Y,Z)$ margins as $P$, and
\[
H(F)-H(P)=I_P(X;Z\mid Y)\ \ge0 .
\]
\end{proposition}
\begin{proof}
$H(F)=H(X,Y)+H(Y,Z)-H(Y)$, and $H(P)=H(X,Y)+H(Y,Z)-H(Y)-I_P(X;Z\mid Y)$.
\end{proof}

So the filler maximizes entropy in its fibre for every $P$. What can be tested on data is not that, but the size of the gap.

%=====================================================================
\section{The update monoid}\label{sec:monoid}
%=====================================================================

An \emph{update} is a partial map from bands to bands: it sends a band $B$ to the closed convex hull of the images of its
members, and it is defined when it is defined on every member. We consider four kinds: conditioning on an event $E$,
$p\mapsto p(\cdot\mid E)$, defined when $p(E)>0$ throughout $B$; the Jeffrey update on a partition $\{E_i\}$ to target weights
$q$, $p\mapsto\sum_iq_ip(\cdot\mid E_i)$ \cite{DiaconisZabell1982}, defined when every $p(E_i)>0$; discounting by reliability
$\alpha$, which sends $p$ to the band of \cref{prop:discount}; and pushforward by a kernel. The hull is needed because the image
of a convex band under a Jeffrey update need not be convex. Composition makes the updates a monoid of partial maps.

\begin{proposition}[Properties of the monoid]\label{prop:monoid}
\begin{enumerate}[label=(\roman*)]
\item Conditioning on events commutes, $p(\cdot\mid E)(\cdot\mid E')=p(\cdot\mid E\cap E')$, and is idempotent.
\item Jeffrey updates on two partitions do not commute in general. They commute at a law under which the two partitions are
  independent.
\item Discounting sends an elementwise band of width $w_y$ at $y$ to one of width $\alpha w_y+1-\alpha$. A kernel $K$ contracts the total-variation diameter of every
  band by at least its Dobrushin coefficient $\eta(K)=\max_{y,y'}\TV(K_y,K_{y'})$.
\item The invertible kernels are the permutations of $Y$. In the monoid $\operatorname{End}(f)$ of kernels that preserve $f$, the
  invertible elements are the permutations $\sigma$ with $f\circ\sigma=f$.
\end{enumerate}
\end{proposition}
\begin{proof}
(i) is direct. (ii) For the $2\times2$ law with rows $(0.30,0.10)$ and $(0.15,0.45)$, updating rows to $(0.6,0.4)$ and columns
to $(0.35,0.65)$ in the two orders gives laws at total variation $0.120$ (check B9). If the partitions are independent, the law
of the two partitions is preserved by both updates, and for $y\in E_i\cap F_j$, with all $p(E_i),p(F_j)>0$, either order gives
$q_ir_j\,p(y)/(p(E_i)p(F_j))$, which is symmetric in the two updates. (iii) The
first statement follows from \cref{prop:discount}, applied to the extreme points. For the second, $\TV(K^\T p,K^\T q)\le\eta(K)\TV(p,q)$ \cite{Dobrushin1956}, as in
\cite[Prop.~11.2]{Tower}. (iv) A stochastic matrix with a stochastic inverse is a permutation matrix, and it preserves $f$ exactly
when $f\circ\sigma=f$; $\operatorname{End}(f)$ is the monoid of \cite[Prop.~3.1]{Algebras}.
\end{proof}

\begin{proposition}[Dilation]\label{prop:dilation}
Let $B$ be the band of laws of $(X,Y)\in\{0,1\}^2$ with $P(X=1)=P(Y=1)=\frac12$. The band for $P(X=1)$ is the point $\frac12$,
but after conditioning on $Y=y$ it is $[0,1]$, for both values of $y$.
\end{proposition}
\begin{proof}
$P(X=1\mid Y=1)=2p_{11}$ with $p_{11}\in[0,\frac12]$ free, and likewise for $Y=0$.
\end{proof}

This is the phenomenon of Seidenfeld and Wasserman \cite{SeidenfeldWasserman1993}: learning something can make every
possible conclusion less certain. It is how unreliable information enters the formalism.

\begin{proposition}[Loops of Jeffrey updates]\label{prop:loop}
On positive two-way tables, Jeffrey updates on the row partition and on the column partition preserve every odds ratio
$p_{ij}p_{kl}/(p_{il}p_{kj})$. A loop that moves the row margin to $q_1$, the column margin to $r_1$, and then back to the
original $q_0$ and $r_0$ in general does not return to the starting law. Iterating the last two updates converges back to it.
\end{proposition}
\begin{proof}
A row update multiplies each row by a constant and a column update each column, and both leave cross ratios unchanged. For
the law in the proof of \cref{prop:monoid}, the loop ends at total variation $0.073$ from the start, with the odds ratio $9$ kept.
Alternating the two updates is iterative proportional fitting, which converges to the unique law with the target margins and the
given odds ratios, the information projection of Csiszár \cite{Csiszar1975}; the starting law is that law.
\end{proof}

So the residue of a loop persists if nothing more is done, but it lies entirely in the margins: the updates leave the interaction
alone, and refitting the original margins removes it. A genuine holonomy, a loop
that changes an invariant, needs updates that act on interactions, such as discounting towards a reference law, and we leave
its study open.

\begin{remark}[Not a Galois action]\label{rem:galois}
The motion of information is not a Galois action, because updates cannot be undone: conditioning forgets and discounting
widens. A Galois-type structure does appear in two places. The units of the update monoid are the relabellings of $Y$, and those that preserve $f$,
the units of $\operatorname{End}(f)$, act on the structures over $f$ and preserve every quantity in this paper; their subgroups
correspond to the partitions they preserve by a Galois connection. And the fibre-level counterpart of a Galois action, monodromy, is the loop question of
\cref{prop:loop}.
\end{remark}

\begin{definition}\label{def:teartime}
For a sequence of margin pairs in which each update moves the pair by at most $v$ in the distance of \cref{def:margins}, the
\emph{tear time} is $\tau=m_T/v$, and the \emph{break time} is $m_B/v$.
\end{definition}

\begin{proposition}[Tear time]\label{prop:teartime}
No structure of \cref{thm:regimes} tears before $\lceil m_T/v\rceil$ updates, and none breaks before $\lceil m_B/v\rceil$ updates.
\end{proposition}
\begin{proof}
By \cref{thm:distances} the walls are at distance $m_T$ and the boundary at distance $m_B$, and each update covers at most $v$.
\end{proof}

The tear time is only a lower bound, because updates need not head for the nearest wall. Whether a ratio of margin to speed is
a useful predictor is an empirical question, and \cref{sec:fhub} gives a mixed answer for a one-dimensional analogue.

%=====================================================================
\section{Families, limits and shards}\label{sec:families}
%=====================================================================

The paper so far moves one law, or one pair of margins, at a time. The structures of \cite{Bridges} live on fibres: for marginal
data $m$, the fibre $C_m=\{\pi\in\D(X):M\pi=m\}$ of compatible joint laws, its maximum-entropy point and its toric and brane data.
This section assembles the fibres into families, says what limits and colimits mean for them, and cuts the parameter space into
regimes.

\subsection{The family over a band}

\begin{proposition}[Families over bands]\label{prop:family}
Let $B$ be a band of marginal data.
\begin{enumerate}[label=(\roman*)]
\item The family $\mathcal P_B=\bigcup_{m\in B}C_m$ equals $\{\pi\in\D(X):M\pi\in B\}$, the preimage of $B$; it is a polytope when $B$ is.
\item For couplings whose margins lie in elementwise bands $F$ and $G$, the upper entropy of $\mathcal P_B$ is $\Hup(F)+\Hup(G)$, attained at
  the product of the two water-fillings of \cref{prop:upper}; the lower entropy is attained at a vertex.
\end{enumerate}
\end{proposition}
\begin{proof}
(i) is the definition of a preimage, and the preimage of a polytope under a linear map, intersected with the simplex, is a polytope.
(ii) $H(\alpha)\le H(f)+H(g)$ for any coupling of $f$ and $g$, with equality at the product, and the product of the water-fillings lies
in $\mathcal P_B$. The lower entropy is a concave minimum over a polytope.
\end{proof}

So the family is not merely a union of polytopes: over a polyhedral band, for example an elementwise one, it is one polytope, and
over any convex band it is convex; its entropy interval is computable. In the lattice of subsets of $\D(X)$, a category different from
the one used for limits below, $\mathcal P_B$ is the colimit, the union, of the fibres $C_m$ over $m\in B$: assembly of the fibres.

\subsection{Limits as compatibility}

Let local laws $m_S$ on $X_S$ be given for the sets $S$ of a cover, agreeing on overlaps. A \emph{cone} over this diagram, in the
category of finite probability spaces and measure-preserving maps, is a probability space with measure-preserving maps to the local spaces, commuting with the marginal maps. Every cone pushes
forward to one with apex $X$ and the coordinate projections as legs, and those are exactly the joint laws $\pi$ with $\pi_S=m_S$: the
fibre $C_m$. The category has no products, and in general there is no universal cone, so ``the limit exists'' has
to be read as ``a cone exists'': compatibility, or a global section in the sense of Abramsky and Brandenburger
\cite{AbramskyBrandenburger2011}. Two refinements make it quantitative. Positive cones exist iff $C_m^\circ\ne\emptyset$, and then entropy
selects a canonical cone, the maximum-entropy law, which is where the entropy brane meets the conormal brane of the constraint
\cite[Thm.~3.1]{Bridges}. When no cone exists, the contextual fraction \cite{AbramskyBarbosaMansfield2017} measures the failure as one
minus the largest weight $\lambda$ of a sub-normalized global section, $\lambda\pi_S\le m_S$ for all $S$.

\subsection{Shards}

\begin{definition}[Structural shards]\label{def:shards}
Let $\Theta$ be a space of parameters $\theta$: marginal data $m$, possibly a pair of margins $(f,g)$ or a band of them, and
possibly a kernel. Its \emph{signature} $\Sigma(\theta)$ has up to three components, each defined when the parameter has the relevant part:
\begin{itemize}[leftmargin=*,nosep]
\item $\sigma_{\mathrm{pos}}$: whether compatible laws exist and, if so, the support of a law in the relative interior of $C_m$ (the
  smallest face of the simplex containing $C_m$); for a band, the set of such supports over its members;
\item $\sigma_{\mathrm{wall}}$, for pairs of margins: the vector of signs of the wall functions $f(I)-g(J)$, with $0$ on a wall; for bands,
  each entry is above, below or straddling, as in \cref{prop:bandsafe};
\item $\sigma_{\mathrm{comp}}$, for kernels: deterministic, constant but not deterministic, or neither.
\end{itemize} A \emph{structural shard}
is a connected component of a level set of $\Sigma$. The \emph{interaction nerve} has the shards as vertices and a simplex for every set
of shards whose closures share a point.
\end{definition}

The name is chosen for its picture, and the notion is unrelated to the shards of hyperplane arrangements \cite{Reading2011}. The
stratification is multiaxial: the three components of $\Sigma$ change independently, so a parameter can sit in the interior of a chamber
and still carry a stochastic kernel. Because $\sigma_{\mathrm{pos}}$ records the support, not only whether a positive law exists, the
dimension of the fibre and of the gluing algebra is constant on a shard. On a shard the structures of \cref{thm:regimes} bend; between shards they tear or break; and
the kernel class decides which compositions are geometric \cite[Thm.~5.1]{Bridges}.

\begin{proposition}[Which updates move between shards]\label{prop:updateregimes}
\begin{enumerate}[label=(\roman*)]
\item Conditioning on an event $E$ with $p(E^c)>0$ breaks: it sets the probabilities outside $E$ to $0$.
\item A Jeffrey update keeps the support of the law if all its targets are positive, and breaks if some target is $0$.
\item Shafer discounting, which sends $p$ to the band $B(\alpha p,\alpha p+1-\alpha)$ of \cref{prop:discount}, keeps $\alpha p$ as the lower
  envelope, so it neither breaks nor heals the lower envelope. Linear discounting towards the uniform law, $p\mapsto\alpha p+(1-\alpha)u$, heals:
  afterwards every probability is at least $(1-\alpha)/n$.
\item A kernel sends a full-support law to a full-support law iff no column of the kernel is zero, and sends every law to a full-support
  law if all its entries are positive.
\item Jeffrey updates, discounting and kernels can cross walls.
\end{enumerate}
\end{proposition}
\begin{proof}
(i)--(iv) are direct; in (iv), $(K^\T p)_j=\sum_yp_yK_{yj}$, and linear discounting is the kernel $\alpha I+(1-\alpha)J/n$, with all entries
positive. For (v), these updates move the margins continuously in their parameters (targets, reliability, kernel entries), and the wall
functions are continuous; for instance a kernel moves a margin $g$ to $K^\T g$, and for suitable $K$ this satisfies $f(I)=(K^\T g)(J)$.
\end{proof}

So among these updates, breaking is created only by conditioning, by Jeffrey updates with a zero target and by kernels with a zero
column, and it is removed by kernels whose entries are positive, linear discounting towards the uniform law among them. Composition
with a kernel that has no zero column can tear but never break, and in the $2\times2$ case it also contracts the moment segment by the
Dobrushin coefficient \cite[Ex.~5.2]{Bridges}. Where the axes interact is visible here: a stochastic composition, an event of the
$\sigma_{\mathrm{comp}}$ axis, moves the margins and can land them on a wall, an event of the $\sigma_{\mathrm{wall}}$ axis.

%=====================================================================
\section{Stability along paths}\label{sec:paths}
%=====================================================================

Two one-parameter families show what bends, what tears and what breaks, structure by structure.

\begin{example}[A $2\times2$ path through two walls]\label{ex:path22}
Fix $g=(0.3,0.7)$ and let $f=(f_1,1-f_1)$ run over $(0,1)$. The walls are $f_1=0.3$ and $f_1=0.7$, and the boundary points are $0$ and $1$.
\begin{center}
\begin{tabular}{p{5.4cm}p{5.6cm}p{3.2cm}}
\toprule
structure & along the path & regime\\
\midrule
gluing algebra $A_f\cong\R\times M_1(\R)$ & constant & bends throughout\\
product coupling $\alpha_{11}=0.3f_1$ & linear & bends throughout\\
moment segment, length $\min(f_1,0.3)-\max(0,f_1-0.7)$ & kinks at $0.3$ and $0.7$, length $\to0$ at the ends & tears, then breaks\\
nonzero torus brane, at the Fréchet midpoint & $f_1/2$, then $0.15$, then $(f_1-0.4)/2$ & tears at both walls\\
vertex entropies & $d\log(1/d)$ cusps at the walls & tear\\
\bottomrule
\end{tabular}
\end{center}
On the walls themselves the margins are degenerate and the toric statements of \cite{Bridges} apply on either side. The product
coupling and the nonzero brane are different points except at $f_1=0.5$ in the middle chamber, and the slope of the brane's
location jumps by $\mp\frac12$ at the two walls. As $f_1\to0$ the toric sphere $X_{f,g}$ shrinks to a point: the geometry breaks with the
margin.
\end{example}

\begin{example}[The three-bit path]\label{ex:path3}
For the pairwise disagreement $1-\varepsilon$ of \cite[Ex.~3.3]{Bridges}, the shards along $\varepsilon\in[0,1]$ are: $[0,\frac13)$, incompatible,
with contextual fraction $1-3\varepsilon$; $\{\frac13\}$, compatible only on the boundary; $(\frac13,1)$, positive; and $\{1\}$, boundary
again. On the positive shard the best minimum cell of a compatible law is $\min(\frac{3\varepsilon-1}4,\frac{1-\varepsilon}4)$, attained by the
maximum-entropy filler, so the filler's break margin falls linearly to $0$ at both ends, while its Fisher form
$\frac8{3\varepsilon-1}+\frac{24}{1-\varepsilon}$ blows up.
\end{example}

%=====================================================================
\section{Second-order effects of breaking}\label{sec:secondorder}
%=====================================================================

A first-order break is a probability reaching $0$. Its second-order effects are what happens, as a consequence, to structures built
on top: global compatibility, the higher levels of the tower, sample sizes, algebras and branes.

\begin{theorem}[Compatibility is lost only through breaking]\label{thm:lossbybreak}
Let $m(s)$ be a continuous path in the affine span of the achievable margins $Q=M(\D(X))$, compatible for $s<s^*$ and incompatible for
$s>s^*$. Then $m(s^*)$ is compatible, every compatible law at $s^*$ lies on the boundary of the simplex, and the largest minimum cell
of a compatible law, $\delta(m)=\max_{\pi\in C_m}\min_x\pi_x$, tends to $0$ as $s\uparrow s^*$.
\end{theorem}
\begin{proof}
$Q$ is a polytope, so it is closed and $m(s^*)\in Q$. If some compatible law at $s^*$ had full support, then $m(s^*)$ would lie in
$M(\D^\circ)$, which is the relative interior of $Q$ because a linear map carries the relative interior of a polytope onto that of its
image \cite[Thm.~6.6]{Rockafellar1970}; the relative interior is open in the affine span, so $m(s)\in Q$ for $s$ slightly above $s^*$, a contradiction. The function
$\delta$ is concave on $Q$, since convex combinations of compatible laws are compatible for the combined margins, and upper
semicontinuous, since $C_m$ depends upper hemicontinuously on $m$; a concave function on a polytope is lower semicontinuous
\cite{GaleKleeRockafellar1968}, so $\delta$ is continuous on $Q$, and $\delta(m(s^*))=0$.
\end{proof}

So a global section cannot disappear while the canonical filler is still strictly positive: the loss of the limit is always announced
by a break of the filler. The next statement turns this into a margin.

\begin{proposition}[Tolerance from the break margin]\label{prop:tolerance}
Let $\pi$ be compatible with $m$ and have minimum cell $\delta$, and let $\kappa=\lVert M^+\rVert_{\infty\to\infty}$ for the pseudo-inverse
$M^+$. Every change $\Delta m$ of the margins within the column space of $M$ that keeps the total mass of every table, with $\lVert\Delta m\rVert_\infty<\delta/\kappa$, leaves
them compatible with a positive law. For pairwise tables of three bits $\kappa=\frac32$, and on the path of \cref{ex:path3}, for $\frac13<\varepsilon\le\frac12$, the bound
$\frac{3\varepsilon-1}6$ equals the distance of the tables to the threshold.
\end{proposition}
\begin{proof}
$\pi+M^+\Delta m$ has margins $m+\Delta m$, total mass $1$ because the all-ones vector lies in the row space of $M$ and every table keeps its
mass, and all cells at least
$\delta-\kappa\lVert\Delta m\rVert_\infty>0$. On the path, $\delta=\frac{3\varepsilon-1}4$ and each table entry moves by $\frac12$ per unit of $\varepsilon$.
\end{proof}

The second-order effects of a break are of five kinds.
\begin{enumerate}[label=(\alph*),leftmargin=*]
\item \emph{Global compatibility} is lost only through a break (\cref{thm:lossbybreak}), with a computable early warning
  (\cref{prop:tolerance}).
\item \emph{Higher levels of the tower.} A zero at level one removes a model at level two for ever, and reliability at level three
  prevents it (\cref{sec:towerdyn}).
\item \emph{Sample sizes.} With joint samples, positivity at distance $m$ from a break is confirmed, by seeing a rare cell, from $O(1/m)$
  samples instead of $O(1/m^2)$, because the Fisher information diverges \cite[\S3]{Bridges}; with pairwise-only samples the rate stays
  $O(1/m^2)$.
\item \emph{Algebra and geometry.} Losing an outcome shrinks the gluing algebra from dimension $(n-1)^2+1$ to $(n-2)^2+1$; its Morita
  type, that of two points, changes only when the last circulation goes ($n=2\to1$), where it becomes that of one point. The toric
  manifold loses dimension, as in \cref{ex:path22}.
\item \emph{A break that does not cascade.} When a middle probability $g_y$ of a chain tends to $0$, the conditional row $K_y$ of the second
  coupling becomes undefined, yet the glued coupling stays continuous: changing $K_y$ moves it by at most $g_y$ in total variation. The break damages the
  kernel, not the composite.
\end{enumerate}

%=====================================================================
\section{The tower under evidence}\label{sec:towerdyn}
%=====================================================================

The tower's calibrated view of an outcome is a level-two law $\Pi=\sum_kw_k\delta_{\rho_k}$ over candidate models, its band
$\operatorname{conv}\{\rho_k\}$ and its barycenter, the predictive law \cite{Tower}. Evidence $x$ updates the weights,
$w_k\propto w_k\rho_k(x)$, which is conditioning at level two, and \cref{prop:level2} splits the entropy of the predictive law at every
step into the expected entropy of the models and the epistemic term $I(K;X)$.

\begin{example}[A break at level one removes a model for ever]\label{ex:cromwell}
Take two models of a three-valued outcome (entropies in nats), $\rho_1=(0.5,0.5,0)$ and $\rho_2=(0.2,0.3,0.5)$, equal prior weights, and the observations
$0,1,0,1,2,0,1,0$. The weight of model 1 and the epistemic term evolve as follows, without and with linear discounting of each model towards the
uniform law with reliability $\alpha=0.9$:
\begin{center}
\begin{tabular}{lcccccccc}
\toprule
observation & $0$ & $1$ & $0$ & $1$ & $2$ & $0$ & $1$ & $0$\\
\midrule
$w_1$, plain & $0.714$ & $0.806$ & $0.912$ & $0.946$ & $0$ & $0$ & $0$ & $0$\\
$I(K;X)$, plain & $0.214$ & $0.186$ & $0.120$ & $0.088$ & $0$ & $0$ & $0$ & $0$\\
$w_1$, discounted & $0.694$ & $0.783$ & $0.891$ & $0.929$ & $0.473$ & $0.671$ & $0.764$ & $0.880$\\
$I(K;X)$, discounted & $0.146$ & $0.125$ & $0.080$ & $0.057$ & $0.152$ & $0.150$ & $0.130$ & $0.085$\\
\bottomrule
\end{tabular}
\end{center}
Without discounting, the first-order break in $\rho_1$ (the zero at outcome $2$) becomes a second-order break of the tower: one
observation removes the model that held $95\%$ of the weight, the band collapses to a point, the epistemic term drops to $0$, and no
later evidence can bring the model back, however strongly it favours it. This is Cromwell's rule seen as breaking. With discounting, the
surprise halves the weight and raises the epistemic term from $0.057$ to $0.152$, so the calibrated view registers the conflict as
uncertainty, and the weight recovers to $0.88$ as the later observations favour model 1.
\end{example}

\begin{proposition}[Reliability bounds how fast the calibrated view can morph]\label{prop:morph}
If every model is linearly discounted towards the uniform law on $n$ outcomes with reliability $\alpha<1$, then each observation changes the
odds between two models by a factor at most $1+\alpha n/(1-\alpha)\le n/(1-\alpha)$. So no weight reaches $0$ after finitely many observations,
and the band of the calibrated view, the hull of the discounted models, stays at distance at least $(1-\alpha)/n$ from every face of the
simplex: it cannot break. Weights may still tend to $0$ as evidence accumulates.
\end{proposition}
\begin{proof}
Discounted probabilities lie between $(1-\alpha)/n$ and $\alpha+(1-\alpha)/n$, and the odds change by the ratio of two of them.
\end{proof}

Read through the tower, the three regimes of the paper become three behaviours of the calibrated view: it bends while evidence moves
the weights, it tears when the predictive law crosses a decision wall, and it breaks when a level-one zero removes a model. Reliability at
level three converts breaking into bending, at the price of a bounded speed.

%=====================================================================
\section{Evidence I: two pilots on public data}\label{sec:pilots}
%=====================================================================

\subsection{Subgroups as contexts: NHANES}\label{sec:nhanes}

\paragraph{Data.} NHANES 2009--10 and 2011--12, as distributed in the R package \texttt{NHANES} \cite{NHANESpkg}: $10{,}730$ adults
aged $20$ or more with blood pressure, body mass index and diabetes status, with examination weights and a stratified jackknife
over primary sampling units for standard errors. Four attributes define $16$ subgroups: age $55$ or more, diabetes, overweight
($\mathrm{BMI}\ge25$) and male sex. The outcome is high systolic pressure ($\ge140$ mmHg), with weighted prevalence $12.8\%$.

\paragraph{The information ladder.} Suppose the sizes of all $16$ subgroups and the overall prevalence are known, and the outcome is
known otherwise only through tables of the outcome with $k$ of the attributes. The risk of a subgroup then lies in a band, computed by linear-fractional
programming over all joints with these tables. For the subgroup "55+, diabetic, overweight, male" ($462$ adults, risk $0.308$):

\begin{center}
\begin{tabular}{lcc}
\toprule
information & band for the subgroup & mean width, all $16$ subgroups\\
\midrule
sizes and overall prevalence & $[0,\ 1]$ & $0.95$\\
outcome by one attribute at a time & $[0,\ 0.881]$ & $0.73$\\
outcome by pairs of attributes & $[0.170,\ 0.430]$ & $0.33$\\
outcome by triples of attributes & $[0.302,\ 0.311]$ & $0.04$\\
\bottomrule
\end{tabular}
\end{center}

This is the precision ladder of \cite[Thm.~3.5]{Attribution} on real data. Every band contains the subgroup's risk, but that holds by
construction, since the bands are computed from the same joint table, and is only a sanity check. The test is out of sample:
build the bands from the $2009$--$10$ tables and ask whether the $2011$--$12$ risks fall in them. With one or two attributes at a time
$16$ and $15$ of the $16$ new risks fall inside. With triples none do: those bands are narrower than the sampling variation between
cycles. When each new risk is given its own $95\%$ sampling interval, $12$ of $16$ intervals meet the triple-level bands. A structural
band alone is overconfident once it is narrower than the sampling noise, which is the point of the next paragraph.

\paragraph{Two kinds of imprecision.} With a risk threshold of $0.30$ as the wall, the structural band from the triples lies entirely
above the wall, while the sampling band, $[0.25,0.37]$, straddles it. (Sampling bands are Wilson intervals with the
design-based effective sample size $r(1-r)/\mathrm{se}^2$.) A tolerance measure has to carry both, the structural band
of the information available and the sampling band of the estimate.

\paragraph{Revision.} Between the two survey cycles the subgroup's risk moved from $0.25$ to $0.36$, outside its earlier $95\%$ band.
Of the $16$ subgroups, $11$ landed inside their earlier bands. With the actual standard errors of both cycles, about $11.3$ of
the $15$ subgroups with positive standard errors would be expected inside if nothing had changed, so $11$ is consistent with no
change. Four subgroups changed status relative to the wall, between decisive and straddling, and none crossed it. One subgroup had $7$ members and no cases in $2009$--$10$; its survey standard error is
exactly $0$, and the band collapses to a point. This is the breaking regime in practice: at the boundary of the simplex the
design-based band is undefined, and we fall back to a count-based interval, $[0,0.35]$.

\paragraph{Contexts and incompatibility.} Pairwise tables taken alternately from the two cycles are slightly incompatible, with least
$L^1$ adjustment $0.084$, but less than tables taken alternately from two halves of the $2009$--$10$ cycle ($0.179$ in the
median of $20$ random splits of its sampling units). Halves have half the sample size, which inflates their noise by about $\sqrt2$, and $0.179/\sqrt2\approx0.13$ still
exceeds $0.084$. So the incompatibility is no larger than sampling noise, and for these tables the two cycles behave as one
population (\cref{sec:protocol}). There is no full Simpson reversal among the
attributes, but the effect of being overweight changes sign with age: $+4.5$ points under $55$ (standard error $0.7$) and $-3.3$
points at $55$ or more ($1.8$), clear in the first stratum and borderline in the second.

\subsection{Revised forecasts and revised data: the COVID-19 Forecast Hub}\label{sec:fhub}

\paragraph{Data.} The ensemble forecasts of the US COVID-19 Forecast Hub \cite{Cramer2022}: weekly incident deaths for $56$ states
and territories, forecast dates June 7, 2021 to March 21, 2022 (targets ending by March 26, so that every target has six weeks of
later data revisions), horizons one to four weeks, $9{,}072$ forecasts given by $23$ quantiles. The truth file of the hub's repository was read at every forecast date, so we know what was known when; the final
truth is the last state-level version (May 9, 2022). The band is the central $80\%$ interval, and each state's wall is its median
weekly deaths over the period.

\paragraph{Calibration and movement.} The bands cover the final value $78$--$80\%$ of the time at every horizon. Revising a forecast
for a fixed week, as the horizon shrinks from four weeks to one, widens the band $26\%$ of the time, and $64\%$ of the time the new
band is not contained in the old one: information moves bands as well as shrinking them.

\paragraph{Tear time.} A band is decisive if it lies on one side of the wall, and it tears when its next revision is not on that side.
The analogue of the tear time used here is one-dimensional: $\tau$ is the margin from the band to the wall, in deaths, divided by
the latest movement of the median, in deaths per revision. It is defined for the $1{,}954$ decisive bands that have a later revision,
an earlier one, and a median that moved; the four-week forecasts ($936$) have no earlier revision and $253$ medians did not move.

\begin{center}
\begin{tabular}{lccccc}
\toprule
$\tau$ (weeks) & $<0.5$ & $0.5$--$1$ & $1$--$2$ & $2$--$4$ & $\ge4$\\
\midrule
share that tear at the next revision & $0.33$ & $0.20$ & $0.11$ & $0.10$ & $0.03$\\
\bottomrule
\end{tabular}
\end{center}

The ranking of decisive bands by $\tau$ separates tearing from not tearing with area under the curve $0.78$. The relative margin
alone does slightly better, $0.80$. So the speed term of $\tau$ adds nothing in this pilot, perhaps because a single movement is
a noisy estimate of speed.

\paragraph{Crossing by straddling.} Of all $3{,}143$ decisive bands with a later revision, only $3$ jumped from one side of the wall to the
other. Bands cross walls by passing through a straddling state, the discrete counterpart of \cref{thm:distances}.

\paragraph{Unreliable data.} Of $2{,}408$ state-weeks, $3.2\%$ were later revised by more than $25\%$. One-week-ahead forecasts made when
the latest week was such a week covered the final value $43\%$ of the time ($95\%$ interval $33$--$55\%$), against $79\%$ otherwise
($74$ and $2{,}278$ forecasts). The bands did not account for the possibility that their inputs would be revised. This is the strongest empirical
reason for a reliability level above the band.

\subsection{What the pilots support}

They support the band picture in four ways: bands narrow with information and, combined with sampling bands, predict new data;
decisive bands rarely jump across walls; small margins predict tearing; and unreliable inputs break calibrated bands. They do
not yet support the tear time as more than a margin, and they do not test holonomy or dilation in the strict sense; the $26\%$ of
forecast bands that widen are not dilations, because successive forecasts condition on different information.

%=====================================================================
\section{Evidence II: a protocol for clinical data}\label{sec:protocol}
%=====================================================================

Clinical records, for example MIMIC-IV \cite{MIMICIV}, offer transitions between patient states, measurements made on different
subsets of patients, and timestamps for when each value was taken and when it was entered. The protocol below separates the
statements that are theorems, which can only fail through errors of code or estimation, from those that could be false.

\begin{enumerate}[leftmargin=*]
\item \emph{Conditional reconstruction.} Estimate transition kernels between clinical states and validate them on held-out patients
  and on a later period. This is a prerequisite, not evidence.
\item \emph{Coupling geometry.} The dimension $(n-1)(m-1)$, the circulation space and the vertices of the coupling polytope follow from
  the margins. The falsifiable part is where the observed joint sits: its tear margin, whether it lies on a face (structural
  zeros), and whether patients whose margins lie near a wall are those whose classification changes between periods.
\item \emph{The canonical filler.} By \cref{prop:gluinggap} the filler always maximizes entropy in its fibre. What can fail is the
  claim that the gap $I(X;Z\mid Y)$ is small when $Y$ is the clinically meaningful intermediate state; it is tested by $2N\,I(X;Z\mid Y)$,
  asymptotically $\chi^2$ with $(\lvert X\rvert-1)(\lvert Z\rvert-1)\lvert Y\rvert$ degrees of freedom.
\item \emph{Higher-order compatibility.} If all tables come from the same patients, a common joint exists: the empirical one. So
  incompatibility can only come from tables computed on different patients, which in clinical data mostly means available-case
  analysis of tests ordered for sicker patients. The observed incompatibility is split into sampling (a simulated or bootstrap null), selection
  (the part removed by inverse-probability weighting for being measured) and a residual that points to differences between
  populations. The measure is fixed in advance: the contextual fraction of Abramsky, Barbosa and Mansfield
  \cite{AbramskyBarbosaMansfield2017}, one minus the largest weight of a sub-normalized common joint below the tables, with the
  least $L^1$ adjustment as a second measure. A successful filling does not count against the theory.
\item \emph{The posterior tower.} That kernels preserve coherence and can only lower the secondary obstruction is a theorem
  \cite[Prop.~9.14]{Tower}, \cite[Prop.~11.5]{Conditionals}. The falsifiable claim is that the estimated obstruction, tear margin or
  tear time predicts deterioration or intervention beyond a standard severity score.
\item \emph{Information arrival.} The two timestamps of each measurement reconstruct what was known at each moment, which is the
  test of \cref{sec:fhub} at the level of a patient.
\end{enumerate}

The thresholds that define walls, the outcomes and the incompatibility measure are to be fixed before the data are examined. The
credentialed data may not be sent to online services unless their handling meets the data use agreement \cite{PhysioNetLLM}, so the
analysis is meant to run where the data are held.

\paragraph{Synthetic validation.} Each test was run on synthetic data with a planted answer (checks S1--S3).
\begin{itemize}[leftmargin=*,nosep]
\item \emph{Test 3.} With $N=2000$ and $X\perp Z\mid Y$ the G-test rejects at level $0.05$ in $4.8\%$ of $400$ replicates; with a planted
  effect of $8$ points it rejects in $88\%$ of $200$.
\item \emph{Test 4.} A lab measured with probability $0.9$, $0.55$ or $0.2$ according to a covariate makes the available-case tables
  incompatible: least $L^1$ adjustment $0.439$ and contextual fraction $0.220$, against a sampling null with median $0.021$ and
  $95\%$ quantile $0.040$ when the lab is missing completely at random. Weighting by the inverse of the estimated probability of
  measurement brings the contextual fraction of the same sample to $0.021$, about the $95\%$ quantile ($0.020$) of the contextual
  fraction under the sampling null, and over repeated samples the weighted inconsistency has median $0.015$ and $95\%$ quantile
  $0.042$, the same order as the sampling null. The decomposition recovers the planted selection.
\item \emph{Walls.} For margins with tear margin $0.080$, the chamber is misidentified from $N$ samples in $24\%$, $0\%$ and $0\%$ of
  $300$ replicates at $Nm_T^2=0.5$, $4$ and $16$. The bound of \cite[Prop.~12.2]{Conditionals}, about $Nm_T^2\ge11$ here for $95\%$, is
  conservative.
\end{itemize}

%=====================================================================
\section{Open problems}\label{sec:open}
%=====================================================================

\begin{enumerate}[leftmargin=*]
\item \emph{Holonomy.} Find updates whose loops change an invariant of the law, for example combinations of discounting and Jeffrey
  updates, and compute the holonomy as a function of the loop.
\item \emph{A better tear time.} Estimate the speed from several revisions, and add the reliability of the inputs, to see whether
  $\tau$ then beats the margin in the forecast data.
\item \emph{Persistence over band width.} As a band widens, the set of chambers it meets grows. The resulting filtration of the chamber
  complex is a persistence module; compute its barcodes for the pilots.
\item \emph{Bands of bands.} Carry the three regimes to levels two and three of the tower, where a band is a set of laws on laws and
  reliability is itself uncertain.
\item \emph{Fisher--Rao margins.} Replace total variation by the Fisher--Rao distance, in which the boundary is at finite distance and
  the break margin and the tear margin live in one metric.
\item \emph{Degeneration of the toric family.} Along a path of margins, describe the limit of the toric manifold, its symplectic potential
  and its nonzero branes at a wall and at the boundary, extending \cref{ex:path22}; and decide whether the positive real Lagrangian of
  \cite{Bridges} extends to a compactification that keeps boundary-supported couplings.
\item \emph{The interaction nerve.} Compute the interaction nerve of the shards for $2\times3$ margins with a kernel axis, and find
  examples where a joint transition carries information that the separate axes do not.
\item \emph{The five papers.} Add to each of the other papers a short subsection on its structures under bands.
\end{enumerate}

%=====================================================================
\section{Computational checks}\label{sec:code}
%=====================================================================

\texttt{check\_bands.py} runs seventeen checks and \texttt{check\_dynamics.py} nine more, all passing. The first:
\begin{itemize}[leftmargin=*,nosep]
\item B1--B3: \cref{prop:discount,prop:tighten,prop:frechet}, by linear programming;
\item B4, B13: \cref{thm:distances,prop:bandsafe};
\item B5--B8: \cref{prop:upper,prop:lower,prop:entropyspeed,prop:entropyregimes,prop:level2,prop:gluinggap};
\item B9--B12: \cref{prop:monoid,prop:dilation,prop:loop,prop:teartime};
\item S1--S3: the synthetic validation of \cref{sec:protocol};
\item P1: the quoted pilot numbers against the saved pilot outputs, and the synthetic and example numbers against this run.
\end{itemize}
The second: D1 (\cref{prop:family}); D2 (\cref{prop:updateregimes}); D3 and D4 (\cref{thm:lossbybreak,prop:tolerance}, on random paths
and perturbations); D5 and D6 (\cref{ex:path22,ex:path3}); D7 and D8 (items (e) and (d) of \cref{sec:secondorder}); D9 (\cref{ex:cromwell,prop:morph},
including the table).
The pilots are reproduced by the scripts \texttt{nhanes\_pilot.py}, \texttt{fhub\_fetch.py} and \texttt{fhub\_pilot.py} in the
folder \texttt{pilot}.

\begin{thebibliography}{99}

\bibitem{Algebras}
V.~K. Awasthi, Coupling algebras: gluing, the square-zero cone and the ring of contingency tables. Technical report, Johns
Hopkins University, October 2026. Draft 4.

\bibitem{Attribution}
V.~K. Awasthi, Precise attribution of model interventions: attribution sets, mechanism reach, path defects and gluing.
Technical report, Johns Hopkins University, October 2026. Draft 5.

\bibitem{Bridges}
V.~K. Awasthi, Bridges: from probability and entropy to branes, with worked examples. Technical note, Johns Hopkins University,
October 2026. Draft 1.

\bibitem{Conditionals}
V.~K. Awasthi, Conditionals, couplings and toric geometry: kernels, the coupling nerve, cumulants and a {F}ukaya
programme. Technical report, Johns Hopkins University, October 2026. Revised draft 4.

\bibitem{ThreeLayers}
V.~K. Awasthi, Three layers of realisation: hierarchy recovery, holonomy, information, contextuality and categorical
dynamics. Technical report, Johns Hopkins University, revised October 2026. DOI: 10.13140/RG.2.2.24274.52163.

\bibitem{Tower}
V.~K. Awasthi, The three-level posterior tower: hierarchical {B}ayesian inference and mixed distributive laws.
Technical report, Johns Hopkins University, October 2026. Revised version 10.

\bibitem{AbellanMoral2003}
J.~Abellán and S.~Moral, Maximum of entropy for credal sets. \emph{Int. J. Uncertainty, Fuzziness Knowledge-Based Systems}
\textbf{11} (2003), 587--597.

\bibitem{AbramskyBarbosaMansfield2017}
S.~Abramsky, R.~S. Barbosa and S.~Mansfield, Contextual fraction as a measure of contextuality. \emph{Phys. Rev. Lett.}
\textbf{119} (2017), 050504.

\bibitem{AbramskyBrandenburger2011}
S.~Abramsky and A.~Brandenburger, The sheaf-theoretic structure of non-locality and contextuality. \emph{New J. Phys.} \textbf{13}
(2011), 113036.

\bibitem{Audenaert2007}
K.~M.~R. Audenaert, A sharp continuity estimate for the von {N}eumann entropy. \emph{J. Phys. A} \textbf{40} (2007), 8127--8136.

\bibitem{Cramer2022}
E.~Y. Cramer et al., Evaluation of individual and ensemble probabilistic forecasts of {COVID}-19 mortality in the {U}nited
{S}tates. \emph{Proc. Natl. Acad. Sci. USA} \textbf{119} (2022), e2113561119.

\bibitem{Csiszar1975}
I.~Csiszár, {$I$}-divergence geometry of probability distributions and minimization problems. \emph{Ann. Probab.} \textbf{3}
(1975), 146--158.

\bibitem{deCampos1994}
L.~M. de~Campos, J.~F. Huete and S.~Moral, Probability intervals: a tool for uncertain reasoning. \emph{Int. J. Uncertainty,
Fuzziness Knowledge-Based Systems} \textbf{2} (1994), 167--196.

\bibitem{DeLoeraKim2014}
J.~A. De~Loera and E.~D. Kim, Combinatorics and geometry of transportation polytopes: an update. In \emph{Discrete Geometry and
Algebraic Combinatorics}, Contemp. Math. 625, Amer. Math. Soc., 2014.

\bibitem{DiaconisZabell1982}
P.~Diaconis and S.~L. Zabell, Updating subjective probability. \emph{J. Amer. Statist. Assoc.} \textbf{77} (1982), 822--830.

\bibitem{Dobrushin1956}
R.~L. Dobrushin, Central limit theorem for nonstationary {M}arkov chains. {I}, {II}. \emph{Theory Probab. Appl.} \textbf{1} (1956),
65--80 and 329--383.

\bibitem{GaleKleeRockafellar1968}
D.~Gale, V.~Klee and R.~T. Rockafellar, Convex functions on convex polytopes. \emph{Proc. Amer. Math. Soc.} \textbf{19} (1968), 867--873.

\bibitem{MIMICIV}
A.~E.~W. Johnson et al., {MIMIC-IV}, a freely accessible electronic health record dataset. \emph{Sci. Data} \textbf{10} (2023), 1.

\bibitem{NHANESpkg}
R.~Pruim, \emph{NHANES: Data from the US National Health and Nutrition Examination Study}. R package version 2.1.4.

\bibitem{PhysioNetLLM}
PhysioNet, Use of {MIMIC} data with large language models and online services.
\url{https://physionet.org/news/post/llm-responsible-use/}.

\bibitem{Reading2011}
N.~Reading, Noncrossing partitions and the shard intersection order. \emph{J. Algebraic Combin.} \textbf{33} (2011), 483--530.

\bibitem{Rockafellar1970}
R.~T. Rockafellar, \emph{Convex Analysis}, Princeton University Press, 1970.

\bibitem{SeidenfeldWasserman1993}
T.~Seidenfeld and L.~Wasserman, Dilation for sets of probabilities. \emph{Ann. Statist.} \textbf{21} (1993), 1139--1154.

\bibitem{Shafer1976}
G.~Shafer, \emph{A Mathematical Theory of Evidence}, Princeton University Press, 1976.

\bibitem{Walley1991}
P.~Walley, \emph{Statistical Reasoning with Imprecise Probabilities}, Chapman and Hall, 1991.

\bibitem{Zhang2007}
Z.~Zhang, Estimating mutual information via {K}olmogorov distance. \emph{IEEE Trans. Inform. Theory} \textbf{53} (2007), 3280--3282.

\end{thebibliography}
\end{document}
